\section{Related work}
\label{sec:related}

\paragraph{Multilingual safety of generative models.} Translating harmful
English requests into other languages is an effective attack on generative
models. \citet{yong2023lowresource} raise the rate at which GPT-4 follows
harmful requests from under 1\% to 79\% by translating them into low-resource
languages, and find much lower rates for high-resource languages.
\citet{deng2023multilingualjailbreak} report that low-resource languages
produce about three times as much harmful content as high-resource ones, and
that combining jailbreak instructions with a non-English language increases
it further. \citet{wang2023alllanguages} build XSafety, covering 14 kinds of
safety issue in 10 languages, and find more unsafe responses in every non-English
language they test. These studies measure what a generator \emph{says}. We
measure what a guard \emph{decides}: whether the classifier placed in front of a
generator still recognises the harmful input in another language.

\paragraph{Multilingual guard models.} PolyGuard \citep{kumar2025polyguard}
trains a safety classifier on 1.91M samples in 17 languages and evaluates it on
PolyGuardPrompts, which combines natural multilingual conversations with
human-verified machine translations. Its finding that existing classifiers
lose accuracy outside English is the closest to ours. Our work differs in the
model class and the protocol: we evaluate typed decision models through one
shared interface, on five tasks that include prompt injection and agent
trajectories, and pair every language with the same model's English answers on
identical cases.

\paragraph{Translated benchmarks.} Machine- or human-translating an English
test set is the standard way to build a multilingual benchmark, as in XNLI
\citep{conneau2018xnli}. Translation is not neutral.
\citet{artetxe2020translation} show that translation introduces artifacts that
models are sensitive to, for example lower lexical overlap between an
independently translated premise and hypothesis. \citet{singh2025globalmmlu}
show that translated knowledge questions carry cultural assumptions of the
source, and that model rankings change on culturally sensitive subsets. Safety
data adds a further problem: its meaning often lives in slang, slurs, typos and
deliberate obfuscation, which a translator may normalise, soften or sharpen. We
report such changes as a finding (Section~\ref{sec:examples}) rather than
assume them away.

\paragraph{Naming the language.} \citet{ducel2022benderrule} find that only
about half of ACL papers name the languages they study, which lets
English-only results read as general ones. We report each language separately,
English included.
